\documentclass[11pt,a4paper]{article}
\usepackage[margin=2.3cm]{geometry}
\usepackage{graphicx}
\usepackage{amsmath}
\usepackage{float}
\usepackage[hidelinks]{hyperref}
\usepackage{caption}

% numbers.tex, counts_table.tex, results_table.tex, transfer_table.tex are written by the notebook
\input{numbers.tex}

\title{Classical Machine Learning for Knee Osteoarthritis Severity Grading\\
\large with External Validation on a Second Hospital Dataset}
\author{NAME-1 (REG-1) \and NAME-2 (REG-2)\\[2pt]
CS446 Machine Learning -- Homework 2, Question 5\\
University of Management and Technology (UMT)}
\date{November 2026}

\begin{document}
\maketitle

\begin{abstract}
Knee osteoarthritis (OA) is graded from radiographs with the five-level Kellgren--Lawrence (KL) scale.
We predict the KL grade from knee X-rays using histogram-of-oriented-gradients (HOG) features, PCA and four
classical classifiers (k-NN, Gaussian naive Bayes, an RBF support vector machine with a hard margin, and a softmax
regression written from scratch). Models are trained and validated on the Osteoarthritis Initiative (OAI) based
Kaggle dataset and tested both internally and on an external dataset from a different hospital. The best model
on validation data was \bestModel{}; it reached a balanced accuracy of \bestKagBA{} on the internal test set but only
\bestDigBA{} on the external set, which shows a clear domain shift. A linear hard-margin SVM separates KL0 from KL4
perfectly on the training set.
\end{abstract}

\section{Introduction}
Osteoarthritis of the knee is a common degenerative joint disease. Doctors read a radiograph and assign a
Kellgren--Lawrence grade~\cite{kl1957}: KL0 healthy, KL1 doubtful, KL2 minimal, KL3 moderate and KL4 severe,
based on joint-space narrowing, osteophytes and sclerosis. The grading is subjective; two experts often disagree
on neighbouring grades. An automatic grader can support screening, but it is only useful if it keeps working on
images from other hospitals. In this homework we build classical (non-deep) classifiers and test them on a
second dataset to measure this.

\section{Datasets}
\textbf{Kaggle / OAI dataset.} Knee Osteoarthritis Severity Grading Dataset by Chen~\cite{chen2018}
(Kaggle copy): $224\times224$ knee-joint crops already split into train, validation and test folders with
sub-folders 0--4 for the KL grade. We used it for training, validation and the internal test.\\
\textbf{Digital Knee X-ray dataset.} Gornale and Patravali~\cite{gornale2020}: 1650 radiographs graded by two
medical experts. We used it only as an external test set with the grades of Expert-I as reference.

\begin{table}[H]\centering
\input{counts_table.tex}
\caption{Number of images per KL grade after removing duplicate files.}
\end{table}

\textbf{Pre-processing.} 16-bit images were converted to floating point, only the left knee of bilateral images
was kept, duplicate files were removed, every image was resized to $64\times64$ and its histogram was
equalised. We checked that no patient (file-name prefix) appears in two Kaggle splits. The classes are strongly
imbalanced (KL0 is the largest class, KL4 the smallest), and the two datasets look different
(Figure~\ref{fig:samples}).

\begin{figure}[H]\centering
\includegraphics[width=0.85\textwidth]{figures/samples.png}
\caption{Example images of each KL grade from the Kaggle dataset (top) and the Digital dataset (bottom).}
\label{fig:samples}
\end{figure}

\section{Methods}
\textbf{Features.} For each image we computed HOG features~\cite{dalal2005} (9 orientations, $8\times8$ cells,
$2\times2$ blocks), giving 1764 values. A \texttt{StandardScaler} and PCA with 100 components (keeping
\pcaVar{} of the variance) were fitted on the training set only and then applied to the validation and test sets,
so that no information leaks from them.

\textbf{Classifiers} (scikit-learn~\cite{sklearn} unless stated otherwise):
\begin{itemize}
\item \textbf{k-NN} on the PCA features; $k\in\{1,3,5,9,15,25,41\}$ and the metric (Euclidean or Manhattan) were chosen
by validation balanced accuracy. Selected: $k=\bestK$, \bestMetric{} distance.
\item \textbf{Gaussian naive Bayes} on the PCA features.
\item \textbf{Hard-margin SVM}~\cite{cortes1995}: \texttt{SVC} with $C=10^{10}$ (not tuned) emulates a hard margin.
A linear kernel on the standardised 1764-D HOG vectors was used for KL0 vs KL4, and an RBF kernel
(one-vs-one, \texttt{gamma="scale"}) on the PCA features for the five grades.
\item \textbf{Softmax regression} implemented with NumPy only: stable softmax, mean cross-entropy with an
$L_2$ penalty $\lambda=10^{-3}$ on the weights (not the bias) and batch gradient descent ($\eta=0.1$, 1500 iterations) on
the standardised PCA features.
\end{itemize}

\textbf{Protocol.} Training on the Kaggle training set, hyper-parameters (k, metric) by balanced accuracy on the
validation set, a single evaluation on the Kaggle test set and on the Digital set. Metrics are accuracy, balanced
accuracy, macro-F1 and quadratic weighted kappa (QWK). Balanced accuracy and QWK are used because the classes are
imbalanced and the grades are ordered.

\section{Results}
\begin{table}[H]\centering
\input{results_table.tex}
\caption{Performance of the four models on the Kaggle test set (internal) and on the Digital set (external).}
\label{tab:results}
\end{table}

\begin{figure}[H]\centering
\includegraphics[width=0.48\textwidth]{figures/knn_curves.png}
\includegraphics[width=0.48\textwidth]{figures/softmax_loss.png}
\caption{Left: validation balanced accuracy of k-NN for different k and both metrics. Right: training and
validation loss of the softmax regression written from scratch.}
\label{fig:curves}
\end{figure}

\begin{figure}[H]\centering
\includegraphics[width=0.95\textwidth]{figures/confusion_best.png}
\caption{Confusion matrices of the best model (\bestModel{}) on the Kaggle test set (left) and on the Digital set (right).}
\label{fig:cm}
\end{figure}

\textbf{Hard-margin SVM check.} On the \svmBinNtrain{} KL0/KL4 training images the linear SVM gave
$\min_i y_i f(x_i)=\svmBinMin{}$, i.e.\ the constraints $y_if(x_i)\ge1$ hold up to rounding and the data are linearly
separable in the 1764-D space. It used \svmBinSV{} support vectors, had a margin width $2/\|\mathbf{w}\|=\svmBinMargin{}$,
and obtained a balanced accuracy of \svmBinBAk{} on the Kaggle test set and \svmBinBAd{} on the Digital set (KL0/KL4 images only).
The RBF-SVM for five grades reached a training accuracy of \rbfTrainAcc{} with \rbfSV{} of \nTrain{} training images as support vectors.

\textbf{Agreement between experts and in-domain training.} Cohen's kappa between the two Digital experts is
\kappaExperts{} (quadratic weighted: \kappaExpertsQ{}). The best model trained within the Digital set (70\,\% stratified split) is
compared with the Kaggle-trained model on the same 30\,\% hold-out in Table~\ref{tab:transfer}.

\begin{table}[H]\centering
\input{transfer_table.tex}
\caption{Best model on the same 30\,\% Digital hold-out: trained on Kaggle (external) vs trained within Digital.}
\label{tab:transfer}
\end{table}

\section{Discussion}
% EDIT THIS SECTION after you have read your own numbers. Keep the sentences that are true for your results.
\textbf{Best model.} On validation data \bestModel{} was best. On the internal test set it obtained accuracy \bestKagAcc{},
balanced accuracy \bestKagBA{} and QWK \bestKagQWK{}.

\textbf{Binary vs five-grade task.} KL0 and KL4 are at opposite ends of the scale and look very different, so a
linear hard-margin SVM can separate them; the five-grade problem contains neighbouring grades (for example KL1 vs KL2)
that even experts confuse, so it is much harder. A hard margin on noisy labels forces the RBF-SVM to memorise noise
(training accuracy \rbfTrainAcc{}, many support vectors).

\textbf{External validation.} On the Digital set the best model obtained accuracy \bestDigAcc{}, balanced accuracy
\bestDigBA{} and QWK \bestDigQWK{}. The difference to the internal test set is caused by domain shift: different
X-ray machines, exposure, cropping and patient population change the HOG/PCA features. Training inside the
Digital set gave balanced accuracy \inDomBA{} compared to \transBA{} for the Kaggle-trained model on the same hold-out,
and the agreement between the two experts (kappa \kappaExperts{}) shows how noisy the reference labels themselves are.

\textbf{Limitations.} Only 64$\times$64 images and hand-crafted HOG features were used, which lose fine details such
as small osteophytes; the hyper-parameters were chosen on one validation split; KL4 has very few images; the second
dataset is graded by only two experts; and no deep learning or data augmentation was tried.

\section{Conclusion}
Classical classifiers on HOG and PCA features can grade knee OA to some extent, and the hard-margin SVM
check confirmed that the extreme grades KL0 and KL4 are linearly separable. However, performance drops on images from another
hospital, so external validation, balanced metrics and ordinal metrics such as QWK are necessary before any clinical use.

\begin{thebibliography}{9}
\bibitem{kl1957} J.~H. Kellgren and J.~S. Lawrence, ``Radiological assessment of osteo-arthrosis,''
\textit{Annals of the Rheumatic Diseases}, 16(4):494--502, 1957.
\bibitem{chen2018} P.~Chen, ``Knee Osteoarthritis Severity Grading Dataset,'' Mendeley Data, V1, 2018,
doi:10.17632/56rmx5bjcr.1.
\bibitem{gornale2020} S.~Gornale and P.~Patravali, ``Digital Knee X-ray Images,'' Mendeley Data, V1, 2020,
doi:10.17632/t9ndx37v5h.1.
\bibitem{dalal2005} N.~Dalal and B.~Triggs, ``Histograms of oriented gradients for human detection,''
in \textit{CVPR}, 2005.
\bibitem{cortes1995} C.~Cortes and V.~Vapnik, ``Support-vector networks,'' \textit{Machine Learning}, 20:273--297, 1995.
\bibitem{sklearn} F.~Pedregosa et al., ``Scikit-learn: Machine learning in Python,''
\textit{Journal of Machine Learning Research}, 12:2825--2830, 2011.
\end{thebibliography}
\end{document}
